% !TeX program = xelatex
\documentclass[11pt]{article}

\usepackage[a4paper, margin=1in]{geometry}
\usepackage{amsmath, amssymb, amsthm}
\usepackage{stmaryrd}
\usepackage{listings}
\usepackage{xcolor}
\usepackage{hyperref}
\usepackage{microtype}
\usepackage{fontspec}
\setmonofont{DejaVu Sans Mono}[Scale=0.85]

\sloppy
\setlength{\emergencystretch}{3em}

\newcommand{\btt}[1]{\texttt{\detokenize{#1}}}

% ---- Lean code listing style -------------------------------------------
\lstdefinelanguage{lean}{
  morekeywords={theorem, lemma, def, by, fun, let, in, do, match, with,
    if, then, else, return, import, open, namespace, end, structure,
    inductive, partial, private, noncomputable, syntax, elab, elab_rules,
    macro, macro_rules, tactic, conv, command, example, have, show, intro,
    apply, simp, rw, ring, all_goals, repeat, decide, native_decide,
    fin_cases, norm_num, exact, set_option},
  sensitive=true,
  morecomment=[l]{--},
  morecomment=[s]{/-}{-/},
  morestring=[b]"
}
\lstset{
  inputencoding=utf8,
  language=lean,
  basicstyle=\ttfamily\footnotesize,
  keywordstyle=\color{blue!70!black}\bfseries,
  commentstyle=\color{gray}\itshape,
  stringstyle=\color{red!60!black},
  showstringspaces=false,
  breaklines=true,
  columns=fullflexible,
  frame=single,
  framesep=4pt,
  rulecolor=\color{black!30},
  backgroundcolor=\color{gray!6},
  xleftmargin=6pt,
  xrightmargin=6pt,
  aboveskip=0.8\baselineskip,
  belowskip=0.8\baselineskip
}
% -----------------------------------------------------------------------

\title{Tactics in LeanFlagAlgebras: an Experience Report}
\author{}
\date{}

\begin{document}
\maketitle

\section{Why we built these tools}

Flag algebras are an active method in extremal combinatorics, but in
informal conversations with combinatorialists we picked up the impression
that the output of flag-algebra tools such as flagmatic is not fully
trusted. We are not aware of a publication that addresses this skepticism
head-on, and the skepticism itself is less an outright rejection than a
reasonable hesitation about verification gaps. Its shape is roughly
threefold: SDP solvers compute in floating point and then round to
rationals, with the resulting certificate expected to be checked for
positive semidefiniteness only after the fact and typically by the user;
the solver pipelines themselves have not been independently verified in
any systematic way; and even when a certificate is produced, it is a
large rational matrix that no human can audit by hand. Reconstructing
the entire certificate inside a proof assistant is one partial response
to this state of affairs, and that is the response we pursue here.

\section{Hand-proving Mantel}

As the smallest possible case study, we first attempted Mantel's theorem.
Only graphs on at most three vertices appear in it, so we believed it could
be carried out by hand, without automation. We did in fact finish that
attempt, and the result is preserved in
\texttt{LeanFlagAlgebras/Archive/MantelTheorem/}.

What struck us about that exercise was not the \emph{length} of the code
but its \emph{shape}. As we wrote out flag definitions, density
calculations, and the product table by hand, it became clear that nearly
every lemma had the same skeleton, with only indices varying. For instance,
the lemma stating ``flag $F_i$ contains $K_3$, hence \texttt{flagDensity} is
zero (or not)'' is mechanical for every $F_i$; the lemma resolving the
product $F_i \cdot F_j$ into a linear combination of host flags has a fixed
form for every pair; the unlabeling and downward lemmas for each unit flag
differ only in their index ranges. None of this is the kind of work a human
should spend effort on. The pattern is fixed; what we needed was a way to
write the pattern down once and have it expanded over the relevant index
sets automatically.

When the hand proof was finished, we drew two conclusions. First, Mantel can
be done by hand. Second, theorems such as $C_4$-Tur\'an or the Erd\H{o}s
pentagon push the host flag size to $5$ and quadratic-form dimensions up to
$8$; the Erd\H{o}s pentagon proof alone involves $8^2 + 6^2 + 5^2 = 125$
flag product terms. Spending human time on more of the same mechanical
repetition is not meaningful. At that point our requirement for the tools
we wanted to build was clear: given a flagmatic-style certificate, the
work that should remain in human hands per theorem is \emph{writing down
one quadratic form} and \emph{a few lines of tactic invocations}.

\section{How a flag algebra proof is written in this framework}

This section follows one concrete proof, \texttt{API/C4TuranAPI.lean}
(a 4-vertex density bound for $K_3$-free graphs; the answer is $3/8$),
through every stage. The other main theorems
(\texttt{MantelTheoremAPI}, \texttt{K4freeP4},
\texttt{ErdosPentagonAPI}) follow a similar skeleton; the details (number
of quadratic forms, host sizes, type structure) differ from theorem to
theorem, but the shape of the file is the same.

\subsection{The overall flow}

An API proof file is laid out in four parts:

\begin{enumerate}
  \item Declare the \textbf{SDP certificate matrices} and prove that each is
    positive semidefinite.
  \item Bulk-load the \textbf{external data} needed for the proof
    (flag enumerations, density tables, product tables) via the
    \texttt{load\_*} meta-commands.
  \item Declare the \textbf{types and flag vectors} that the matrices are
    paired with.
  \item State the theorem and discharge it with a \textbf{stereotyped
    sequence of tactics}.
\end{enumerate}

Parts~1 and~3 are where the numerical content of the certificate enters
Lean. Part~2 registers the objects (which flag, which product) that the
numerical content refers to. Part~4 is the proof body itself.

Both Parts~1 and~2, as we will see, are written by tools rather than by
hand. The author of a new theorem supplies matrix entries and a small
amount of metadata as JSON; Python scripts shipped with the repository then
produce the Lean PSD blocks and the JSON tables that the
\texttt{load\_*} commands consume.

\subsection{Part 1: SDP certificate matrices and positive semidefiniteness}

The heart of the flag algebra method is the fact that, for a positive
semidefinite matrix $M$ and a vector $v$ of flags over a type $\sigma$, the
quadratic form $v^\top M v$ can be added to either side of a density
inequality without breaking it. For $C_4$-Tur\'an the certificate consists
of two matrices: $M_1$ ($4 \times 4$, paired with the type $\sigma_1$ of
two non-adjacent vertices) and $M_2$ ($3 \times 3$, paired with $\sigma_2$
of two adjacent vertices). The SDP certificates we work with come already
distilled to exact rationals, so transcribing one of them into Lean looks
like this:

\begin{lstlisting}
def M_1 : Matrix (Fin 4) (Fin 4) ℚ :=
  !![( 3/8 : ℚ),  0, 0, (-3/8 : ℚ);
       0,         0, 0,   0;
       0,         0, 0,   0;
     (-3/8 : ℚ),  0, 0, ( 3/8 : ℚ)]
noncomputable def M_1_real : Matrix (Fin 4) (Fin 4) ℝ :=
  ratMatrixToReal M_1
\end{lstlisting}

We still owe Lean a proof that $M_1$ is positive semidefinite. We discharge
every such obligation in a uniform way, by exhibiting an explicit
$\mathrm{LDL}^\top$ decomposition: we produce an $L$ and a non-negative
diagonal $d$ such that $M = L \cdot \mathrm{diag}(d) \cdot L^\top$, and let
a general lemma turn this into PSD-ness. For each matrix the resulting Lean
block has the following shape:

\begin{lstlisting}
def dM_1 : Fin 4 → ℚ := ![3/8, 0, 0, 0]
def LM_1 : Matrix (Fin 4) (Fin 4) ℚ :=
  !![1, 0, 0, 0;
     0, 1, 0, 0;
     0, 0, 1, 0;
    -1, 0, 0, 1]
lemma dM_1_nonneg (i : Fin 4) : 0 ≤ dM_1 i := by
  fin_cases i <;> norm_num [dM_1]
lemma M_1_eq_LDL : M_1 = LM_1 * Matrix.diagonal dM_1 * LM_1ᵀ := by
  decide +kernel
theorem M_1_posSemidef : M_1.PosSemidef :=
  posSemidef_of_eq_mul_diagonal_mul_transpose dM_1_nonneg M_1_eq_LDL
-- The same decomposition is then lifted over ℝ to obtain M_1_real_posSemidef.
\end{lstlisting}

\paragraph{Producing this block automatically.}
The \emph{shape} of this block is the same for every matrix, but finding
the actual $L$ and $d$ by hand is tedious. For that we provide
\texttt{LeanFlagAlgebras/Utils/Matrix/generate\_psd\_proof.py}. Given the
matrix as a JSON array of rational rows (the file shipped alongside the
script as \texttt{matrix.json}, for instance), the script computes the
exact-rational $\mathrm{LDL}^\top$ decomposition itself and emits the whole
Lean block. The usage is a one-liner:

\begin{lstlisting}[language={}]
python generate_psd_proof.py --input matrix.json --name M_1 \
       --out C4TuranAPI.lean --write-mode append
\end{lstlisting}

This produces \texttt{M\_1}, \texttt{dM\_1}, \texttt{LM\_1},
\texttt{dM\_1\_nonneg}, \texttt{M\_1\_eq\_LDL},
\texttt{M\_1\_posSemidef}, and the corresponding $\mathbb{R}$-level
statements, ready to paste.

\paragraph{What ends up in the trusted base.}
The correctness of the script itself does not. If the script emits a bad
$L$ or $d$, the call to \texttt{decide +kernel} in
\texttt{M\_1\_eq\_LDL} fails, the file does not compile, and the
discrepancy is surfaced. The script merely \emph{searches for} an
$\mathrm{LDL}^\top$ decomposition; whether the proposed decomposition is
really one of $M$ is re-checked by the Lean kernel.

\subsection{Part 2: Bulk-loading external data}

Quite apart from the matrix data, the proof needs a number of other
numerical facts registered in advance: which size-$n$ host flags are
$K_3$-free, what the pair density of each $\sigma_i$-flag pair is on each
size-$n$ host flag, and how each such product expands as a linear
combination of host flags. In this theorem the following five lines take
care of all of these:

\begin{lstlisting}
load_forbid_density_theorems
  "LeanFlagAlgebras/Flags/Densities/graphs_4_K3_free_indices.json"
load_flag_pair_density_theorems
  "LeanFlagAlgebras/Flags/Densities/density_4_2_0_from_3_2_0_forbid_K3.json"
load_forbid_mul_theorems
  "LeanFlagAlgebras/Flags/Densities/density_4_2_0_from_3_2_0_forbid_K3.json"
load_flag_pair_density_theorems
  "LeanFlagAlgebras/Flags/Densities/density_4_2_1_from_3_2_1_forbid_K3.json"
load_forbid_mul_theorems
  "LeanFlagAlgebras/Flags/Densities/density_4_2_1_from_3_2_1_forbid_K3.json"
\end{lstlisting}

What each line registers, in brief (see \texttt{tactics.tex} for the
full specification of each meta-command):

\begin{itemize}
  \item \texttt{load\_forbid\_density\_theorems} registers, as
    \texttt{@[simp]} lemmas, the statement ``flag $F_i$ is (or is not)
    $K_3$-free'' for every size-4 empty-typed flag.
  \item \texttt{load\_flag\_pair\_density\_theorems} registers, as
    \texttt{@[simp]} lemmas, the pair density of every $\sigma_i$-flag pair
    on every size-4 host flag.
  \item \texttt{load\_forbid\_mul\_theorems} pulls the product table from
    the same JSON and registers, for every relevant pair $(F_i, F_j)$, the
    product expansion lemma
    $F_i \cdot F_j =_{K_3\text{-forbid}}
       \sum_h c_{ijh}\, F_h$, named \texttt{flagMul\_<A>\_<B>}.
\end{itemize}

The base definitions referenced by these statements (\texttt{Flag\_*},
\texttt{FlagAlgebra\_*}) are themselves registered project-wide in
\texttt{Flags/FlagDef.lean} via \texttt{load\_empty\_typed\_flags} and
\texttt{load\_flags}, so individual proof files do not have to repeat them.

\paragraph{Producing these JSON files automatically.}
None of the JSON files referenced above are written by hand. They are
produced by Python scripts in \texttt{LeanFlagAlgebras/Flags/} and its
\texttt{Graphs/}, \texttt{Flags/}, and \texttt{Densities/} subfolders:

\begin{itemize}
  \item \texttt{Flags/generate\_graphs.py} takes a size $n$ and emits
    \texttt{Graphs/graphs\_n.json}, the canonical enumeration of
    empty-typed flags of size $n$. This is the input to
    \texttt{load\_empty\_typed\_flags}.
  \item \texttt{Flags/generate\_flags.py} takes $(n, k, m)$ and emits
    \texttt{Flags/flags\_n\_k\_m.json} containing every $\sigma_m$-flag of
    size $n$ together with its type embedding and downward normalizing
    factor. This is the input to \texttt{load\_flags}.
  \item \texttt{Flags/Densities/calculate\_densities.py} takes pattern and
    host indices and emits the pair-density tables
    \texttt{density\_*.json}, consumed by both
    \texttt{load\_flag\_pair\_density\_theorems} and
    \texttt{load\_forbid\_mul\_theorems}.
  \item \texttt{Flags/Densities/gen\_free\_indices.py} takes a forbidden
    subgraph $H$ and emits, for the requested size $n$, the index list of
    $H$-free flags, consumed by
    \texttt{load\_forbid\_density\_theorems}.
\end{itemize}

In other words, supporting a new forbidden subgraph or a larger host size
amounts to running the relevant script once and adding the corresponding
\texttt{load\_*} call to the proof file. As before, none of this places
the scripts in the trusted base: every numerical fact loaded from JSON is
re-checked by \texttt{decide} or \texttt{native\_decide} inside Lean, and a
mis-generated JSON shows up as a compile error.

\subsection{Part 3: Declaring types and flag vectors}

We then name the flag vector $v_i$ that each matrix $M_i$ will be paired
with:

\begin{lstlisting}
def σ₁ : FlagType (Fin 2) := FlagType_2_0
noncomputable def v₁ : FlagAlgebraVec σ₁ 4 :=
  ![FlagAlgebra_3_2_0_0, FlagAlgebra_3_2_0_1,
    FlagAlgebra_3_2_0_2, FlagAlgebra_3_2_0_3]

def σ₂ : FlagType (Fin 2) := FlagType_2_1
noncomputable def v₂ : FlagAlgebraVec σ₂ 3 :=
  ![FlagAlgebra_3_2_1_0, FlagAlgebra_3_2_1_1, FlagAlgebra_3_2_1_2]
\end{lstlisting}

Every name appearing here was registered earlier by \texttt{load\_flags}.
The mapping ``the $i$-th entry of $v_j$ refers to flag $X$'' is mediated
not by an extra lookup table but by the numeric suffixes $(n,k,m,i)$ in
the names themselves.

\subsection{Part 4: The proof body}

Now the theorem statement and proof body.

\begin{lstlisting}
theorem C4_flagAlgebra_API
    : FlagAlgebra_4_0_0_8 ≤[K3.toFinFlag] (3 / 8 : ℝ) • (1 : FlagAlgebra ∅ₜ)
  := by
  have quadraticForm_trans : FlagAlgebra_4_0_0_8 ≤[K3.toFinFlag]
        FlagAlgebra_4_0_0_8
        + ⟦flagQuadraticForm M_1_real v₁⟧₀
        + ⟦flagQuadraticForm M_2_real v₂⟧₀ := by
    apply forbidLE_add_QuadraticForm M_2_real M_2_real_posSemidef v₂
    apply forbidLE_add_QuadraticForm M_1_real M_1_real_posSemidef v₁
    exact forbidLE_refl K3.toFinFlag FlagAlgebra_4_0_0_8
  apply forbidLE_trans quadraticForm_trans
  apply forbidLE_trans_forbidEq_right ?_
        (forbidEq_smul (forbidEq_symm
          (one_forbidEq_forbidExpand_one K3.toFinFlag 4)))
  simp [flagQuadraticForm, v₁, M_1_real, ratMatrixToReal, M_1,
        Fin.sum_univ_four, add_assoc]
  simp [v₂, M_2_real, ratMatrixToReal, M_2,
        Fin.sum_univ_three, add_assoc]
  reduce_downward_flagmul
  expand_one_at 4
  simp [smul_smul, downward_add, downward_smul]
  ac_sort_rhs_pipeline
  apply forbidLE_of_le
  flag_nonneg
\end{lstlisting}

We walk through what each line does.

\paragraph{\texttt{have quadraticForm\_trans : \ldots}}
Since each $M_i$ is positive semidefinite, the quadratic form
$\llbracket v_i^\top M_i v_i \rrbracket_0$ is non-negative under the
forbid condition, so adding it to the left-hand side of an inequality
preserves the inequality. Applying this twice extends our left-hand side
to ``original term $+$ first quadratic form $+$ second quadratic form.''
The helper lemma \texttt{forbidLE\_add\_QuadraticForm} accepts the matrix,
its PSD proof, and the flag vector and performs each such step.

\paragraph{\texttt{apply forbidLE\_trans quadraticForm\_trans}}
Chains the auxiliary inequality just built into the goal's left-hand side.
The goal becomes ``extended left-hand side
$\le_{\text{forbid}}$ original right-hand side.''

\paragraph{\texttt{apply forbidLE\_trans\_forbidEq\_right \ldots\,
one\_forbidEq\_forbidExpand\_one \ldots}}
Replaces the $1 \in \mathrm{FlagAlgebra}$ on the right-hand side by the
weighted sum of size-4 $K_3$-free flags it is equal to under the forbid
condition. The actual expansion happens later, at \texttt{expand\_one\_at
4}; this step only secures the right to make that replacement.

\paragraph{The two \texttt{simp [flagQuadraticForm, \ldots]} lines}
Unfold each $\llbracket v_i^\top M_i v_i \rrbracket_0$ into an explicit sum
of products. \texttt{Fin.sum\_univ\_four} and
\texttt{Fin.sum\_univ\_three} unfold the index sums $\sum_{i<4}$ and
$\sum_{i<3}$, and \texttt{add\_assoc} right-associates the result. After
these two lines the left-hand side is a flat sum of
``coefficient $\cdot (A \cdot B)$'' product terms.

\paragraph{\texttt{reduce\_downward\_flagmul}}
Resolves each product $A \cdot B$ using the
\texttt{flagMul\_<A>\_<B>} lemma loaded in Part~2, turning it into a
linear combination of base-typed flags; forbidden host flags are dropped by
\texttt{simp} in the process. A single invocation of this tactic dispatches
all $4^2 + 3^2 = 25$ product terms in $C_4$-Tur\'an at once.

\paragraph{\texttt{expand\_one\_at 4}}
Now actually rewrite the $1$ on the right-hand side, whose replacement
right was reserved earlier, as a weighted sum of size-4 $K_3$-free flags.

\paragraph{\texttt{simp [smul\_smul, downward\_add, downward\_smul]}}
Distributes scalar multiplication and \texttt{downward} so that every term
takes the flat shape ``coefficient $\cdot$ \texttt{Flag\_4\_0\_0\_*}''.

\paragraph{\texttt{ac\_sort\_rhs\_pipeline}}
Sorts the right-hand side by base-flag index and merges coefficients of
terms sharing the same base. The two sides now match term by term.

\paragraph{\texttt{apply forbidLE\_of\_le}}
Drops the forbid qualifier and reduces the goal to a plain $\le$ inequality
in the flag algebra.

\paragraph{\texttt{flag\_nonneg}}
The remaining goal is of the form ``$0 \le$ non-negative linear combination
of unit flags,'' which this tactic closes in one line.

\end{document}
